\begin{frame}{교훈: 경로 의존(path dependence)}
  \begin{block}{협력 게임에서 학습의 문제는}
    ``최해를 \textbf{찾느냐}''가 아니라 \textbf{``어느 균형으로
    수렴하느냐''} --- \textbf{경로 의존(path
    dependence)}. 초기 탐색의 \textbf{잡음}(어느 쪽에 먼저
    방문 횟수가 쌓였느냐)이 최종 균형을 결정.
  \end{block}
  \vskip0.6em
  \begin{itemize}
    \item 현실에서도 이 구조는 \textbf{흔하다}:
    \begin{itemize}
      \item \textbf{VHS와 Betamax}의 비디오테이프 포맷
        전쟁(1980년대) --- 기술적 우열이 아니라 ``어떤
        쪽이 \textbf{먼저 잠금}되는지''를 가른 조율 게임
      \item 오늘날의 \textbf{통신 표준}(Wi-Fi, USB-C) 경쟁도
        같은 구조 --- 표준화 기구가 존재하는 이유도
        ``학습으로 균형이 \textbf{잘못 잠기면 안 된다}''는
        인식
    \end{itemize}
  \end{itemize}
\end{frame}
